\documentclass[../../../include/open-logic-section]{subfiles}

\begin{document}

\olfileid{sth}{replacement}{strength}
\olsection{替换公理模式的力量}

关于替换公理模式的力量，我们先看一个简单的观察：除非超出 $\Z$，否则我们无法证明任何大于等于 $\omega + \omega$ 的 von Neumann（冯·诺伊曼）序数的存在。

以下是原因的简要说明。在~$\ZF$ 中，考虑集合 $V_{\omega+\omega}$。这个集合可以作为~$\Z$ 的一个\emph{模型}的论域。为了说明这一点，我们先为公式的\emph{相对化}引入一些记号：
\begin{defn}\ollabel{formularelativization}
	对于任意集合 $M$ 和任意公式 $\phi$，令 $\phi^M$ 表示将 $\phi$ 中所有量词限制到 $M$ 上所得出的公式。也就是说，把“$\lexists[x]$”替换为“$(\lexists[x \in M])$”，并把“$\lforall[x]$”替换为“$(\lforall[x \in M])$”。
\end{defn}\noindent
可以证明，对于~$\Z$ 的每个公理 $\phi$，都有 $\ZF \vdash
\phi^{V_{\omega+\omega}}$。但是，根据 \olref[spine][rank]{ordsetrankalpha}，$\omega+\omega$ 并不\emph{在} $V_{\omega+\omega}$ 中。所以，$\Z$ 与“$\omega+\omega$ 不存在”是一致的。

这就是为什么我们在 \olref[ordinals][replacement]{sec} 中说，没有替换公理模式就无法证明 \olref[ordinals][ordtype]{thmOrdinalRepresentation}。因为很容易在~$\ZF$ 内定义一个显式的良序，它在直觉上\emph{应该}具有序型 $\omega+\omega$。事实上，我们在 \olref[ordinals][idea]{sec} 中已经给出了一个非正式的例子，当时我们给出了自然数上的如下排序：
\begin{align*}
	n \lessdot m \text{ 当且仅当 }&\text{要么 }n < m\text{ 且 }m-n\text{ 是偶数，}\\
	& \text{要么 $n$ 是偶数且 $m$ 是奇数。}
\end{align*}
但如果 $\omega+\omega$ 不存在，这个良序就不与任何序数同构。因此，$\Z$ \emph{不能}证明 \olref[ordinals][ordtype]{thmOrdinalRepresentation}。

反过来说：替换公理模式使我们能够证明 $\omega+\omega$ 的存在，因此也必然使我们能够证明 $V_{\omega+\omega}$ 的存在。不仅如此。对于我们可以定义的\emph{任何}良序，\olref[ordinals][ordtype]{thmOrdinalRepresentation} 都告诉我们，存在某个与之同构的 $\alpha$，从而 $V_\alpha$ 存在。于是，替换公理模式以一种直接的方式保证了集合层级必定\emph{非常高}。

在接下来的几节以及 \olref[card-arithmetic][fix]{sec} 中，我们将更好地体会替换公理模式究竟迫使集合层级变得\emph{多}高。目前简单的要点是：替换公理模式确实\emph{需要}证立！

\end{document}